\section{Recovering a unitary representation from its compact types}

% BEGIN M90H1
\paragraph{M90H1. The compact subgroup and its weight lattice.}
For $n=p+q>0$ the group specified by the two homomorphisms in
$\mathrm{GrpData}$ is
\[
 \widetilde G=\{(g,z)\in U(p,q)\times{\mathbb C}^{\times}:z^2=\det g\}.
\]
The specified bijection and embedding give a topological group
isomorphism. The prescribed path above $\exp(tX)$ starts at the
identity and hence equals the unique exponential lift. Put
$K=U(p)\times U(q)$ and let $\widetilde K$ be its inverse image.
This is a maximal compact subgroup: its projection is maximal
compact and it contains the full kernel. Both $\widetilde K$ and
$\widetilde G$ are connected. Indeed $K$ and $U(p,q)$ are connected,
and the loop rotating a single diagonal coordinate once lifts from
$(1,1)$ to $(1,-1)$. The reductive group results below apply to
this linear group: it is a real form of the connected complex
reductive group $\{(a,z)\in GL_n({\mathbb C})\times
{\mathbb C}^{\times}:\det a=z^2\}$. The involution for this real
form is $(a,z)\mapsto(J(a^*)^{-1}J,\bar z^{-1})$ with
$J=\operatorname{diag}(I_p,-I_q)$.

Over the diagonal torus the exponential period lattice, using angles
in units of $2\pi$, consists of the integer vectors whose coordinate
sum is even. Its dual is
\[
   {\mathbb Z}^n+\tfrac12{\mathbb Z}(1,\ldots,1).
\]
Indeed the lattice contains $2e_i$ and $e_i-e_j$, so the dual
coordinates are half-integral with integral pairwise differences;
these conditions also suffice. In particular all torus weights of
a compact type are rational in the matrix coordinates of the source.
% END M90H1

% BEGIN M90H2
\paragraph{M90H2. Compact-finite vectors have analytic Hilbert-space orbits.}
First take an irreducible strongly continuous unitary representation
of the ordinary group $G=U(p,q)$. If $pq>0$, the scalar circle acts
by one character by Schur's lemma. The product of that circle with
$SU(p,q)$ maps onto $G$ by a local analytic isomorphism. Restriction
to $SU(p,q)$ is therefore irreducible, and a $K$-finite vector $v$
is finite under $K\cap SU(p,q)$. Kr\"otz and Stanton,
\emph{Holomorphic extensions of representations. I}, Theorem~3.1,
gives a holomorphic extension of the Hilbert-valued orbit map for
that vector on $SU(p,q)$. Multiplication by the analytic character
of the scalar circle now gives a real analytic $G$-orbit map.
If $pq=0$, an irreducible representation of the compact group is
finite-dimensional and analytic. The trivial group gives the same
conclusion in rank zero. For an irreducible representation of
$\widetilde G$, the central involution acts by $(-1)^j$ with
$j\in\{0,1\}$. Tensoring by the character $(g,z)\mapsto z^{-j}$
makes this representation descend to $G$ and preserves the finite
compact-orbit spans. Pull back the orbit map and restore the analytic
character to obtain the assertion on the cover.

For the ordinary group the source asks for an analytic extension to
an open set of real matrices with complex entries, viewed in their
real and imaginary coordinates. Write $A^\dagger=JA^*J$. Near any
$g_0$ in $G$ the convergent power series at the identity gives
\[
 R(A)=A(A^\dagger A)^{-1/2}.
\]
The inverse square root commutes with $A^\dagger A$ and is fixed by
$\dagger$, since its power series has real coefficients. Thus
$R(A)^\dagger R(A)=I$ and $R(g)=g$ on the group. Composing the
intrinsic analytic orbit map with $R$ proves the precise ambient
extension requested by $\mathrm{KFiniteAnalyticUJ}$. Permuting
coordinates handles every ordering of the signature signs.
% END M90H2

% BEGIN M90H3
\paragraph{M90H3. The compact types form the irreducible admissible module.}
Let $H$ be a nonzero irreducible unitary representation of
$\widetilde G$ and $V$ the algebraic sum of its finite compact
types. Every $v\in V$ is analytic by M90H2. The differentiated
action preserves $V$: if $v$ lies in the finite-dimensional
$\widetilde K$-space $E$, the compact orbit of $d\pi(X)v$ belongs
to $d\pi(\mathfrak g_{\mathbb C})E$, a finite-dimensional space,
by covariance of differentiation. Taking $V$ itself for the
derivative-stable space gives the source's smoothness assertion.
Complexifying the differentiated action gives the matrix Lie algebra
$\mathfrak{gl}_n({\mathbb C})$. Each compact span decomposes into
diagonal weight spaces with the rational weights of M90H1. These
facts give the required abstract module and its inclusion into $H$.
Any other model in the statement maps bijectively onto $V$ and
intertwines the real derivatives, hence also their complexifications.

Harish-Chandra's admissibility theorem gives finite multiplicity of
each compact type (Borel, \emph{Automorphic forms on reductive
groups}, \S8.5). For the chosen block upper triangular positive
roots, the highest-weight theorem for compact connected groups says
that one copy of the type of highest weight $\mu$ contributes one
line killed by the positive compact root operators in weight $\mu$.
All other types contribute zero. Thus the dimension of the source's
highest-vector space is exactly the desired multiplicity, including
zero for weights that do not occur.

The underlying $(\mathfrak g_{\mathbb C},\widetilde K)$-module
of an irreducible unitary representation is irreducible
(Adams, van Leeuwen, Trapa and Vogan, \emph{Unitary representations
of real reductive groups}, Theorem~2.7(1)--(2)).
Here this gives irreducibility even for the Lie algebra alone.
If $N\subset V$ is Lie-invariant, its intersection with each
finite-dimensional compact span is invariant under
$\operatorname{Lie}(\widetilde K)$ and hence under the connected
group $\widetilde K$. This makes $N$ a submodule for the pair.
% END M90H3

% BEGIN M90H4
\paragraph{M90H4. Infinitesimal equivalence and the converse.}
Suppose that the two modules of irreducible unitary representations
are isomorphic by a linear map commuting with the matrix Lie algebra.
M90H3 identifies the map with an isomorphism of the compact-finite
vectors. For a finite-dimensional compact span $E$ in the domain,
the images of a basis lie in a common finite-dimensional compact
span $F$ in the target. The map $E\to F$ intertwines the compact Lie
algebra, so it intertwines the exponentials and hence the connected
group $\widetilde K$. The two irreducible unitary representations
are now infinitesimally equivalent in the usual sense. Harish-Chandra's
unitary equivalence theorem gives a unitary intertwiner of the Hilbert
spaces (Adams--van Leeuwen--Trapa--Vogan, Lemma~2.4 and
Theorem~2.7(3)). Its existence is exactly the conclusion requested.

Conversely, let an arbitrary unitary representation have its entire
compact-finite space $V$ identified with an algebraically irreducible
module as in the statement. Then $V$ is nonzero and dense by
Peter--Weyl. If $L\subset H$ is closed and invariant, its orthogonal
projection $P$ commutes with the group and sends $V$ to $V$.
It commutes on $V$ with each real derivative by boundedness and thus
with the complex Lie algebra. Both $P V$ and $(1-P)V$ are submodules,
dense in $L$ and $L^\perp$, respectively. Algebraic irreducibility
forces one of them to vanish, so $L=0$ or $L=H$. This proves the
converse in the stated generality. Together M90H1--M90H4 establish
all three module clauses and the analytic-orbit clause of
$\mathrm{Published.HarishChandraModules}$.
% END M90H4
